\documentclass[11pt]{article}
\usepackage[margin=25mm]{geometry}
\usepackage{fontspec}\setmainfont{Latin Modern Roman}
\usepackage{amsmath,amssymb,mathtools}
\usepackage{unicode-math}\setmathfont{Latin Modern Math}
\usepackage{xurl,hyperref}
\hypersetup{colorlinks=true,urlcolor=blue}
\setcounter{secnumdepth}{0}\setlength{\parindent}{0pt}\setlength{\parskip}{5pt}\setlength{\emergencystretch}{3em}
\newcommand{\sref}[2]{\hyperlink{source:#1:#2}{[S#2]}}
\newcommand{\cataloguescope}[1]{\paragraph{Recorded target.}#1}
\begin{document}
\section{Kourovka Notebook Problem 21.87}\label{statement}

\subsection{Definitions and mathematical statement}

\paragraph{Shared notation.}$\mathbb N=\{1,2,\ldots\}$, and $\mathbb Z,\mathbb Q,\mathbb R,\mathbb C$ denote the integers, rationals, reals and complex numbers. Any other conventions are specified locally.


For a finite group $G$, write $d(G)$ for its minimal number of generators, and $[G:H]=|G|/|H|$ for a subgroup index. A family of indices has no common divisor when its greatest common divisor is $1$. All subgroups are ordinary abstract subgroups and no normality is required.
\paragraph{Statement recorded in the supplied source.}
For every finite group $G$ and integer $d\geq0$, if there are subgroups $H_1,\ldots,H_m\leq G$ with $d(H_i)\leq d$ for every $i$ and
\[
 \gcd([G:H_1],\ldots,[G:H_m])=1,
\]
must $d(G)\leq d+1$? \sref{statement}{2}
\subsection{Short English statement}
If a finite group has d-generator subgroups with relatively prime indices collectively, must the whole group have at most d+1 generators?
\subsection{Sources}
\begin{description}
\item[\hypertarget{source:statement:1}{S1}] ULAM.AI and the UnsolvedMath Contributors, \emph{UnsolvedMath: A Curated Collection of Open Mathematics Problems}, catalogue entry. CC BY 4.0. \url{https://huggingface.co/datasets/ulamai/UnsolvedMath}.
\item[\hypertarget{source:statement:2}{S2}] A. Lucchini, Problem 21.87, in E. I. Khukhro and V. D. Mazurov (eds.), \emph{The Kourovka Notebook: Unsolved Problems in Group Theory}, 21st edition, arXiv:1401.0300v48 (6 October 2026). This version records an affirmative answer; the displayed question is the historically registered target. \url{https://arxiv.org/abs/1401.0300v48}.
\end{description}
\label{end:statement}
\end{document}
